\section{Introduction}

Digital misinformation often reuses genuine media outside its original context. A photograph from one event may be paired with a claim about another place, date, or incident. Pixels alone may show no manipulation, and image--text compatibility can identify a mismatch without establishing what the image originally depicted. NewsCLIPpings and VERITE demonstrate both the importance of false-context detection and the danger of unimodal shortcuts \cite{luo2021newsclippings,papadopoulos2024verite}. A fact-checking system therefore needs external evidence that connects the claim and image to documented events.

Retrieval supplies that context, but a valid historical decision depends on when the evidence became observable and where it originated. For a claim assessed at cutoff time $\tau$, a later correction, official report, or fact-check page contains information that was unavailable at the decision point. A self-reported publication field is not enough to prove availability: a page first captured after $\tau$ may have been backdated, changed, or unavailable to the contemporary system. Search results can also give a false impression of corroboration. Syndicated articles, reposts, and quotation chains may all trace to one report, so document-level aggregation can overweight one underlying source.

Evidence availability is different from model confidence. The task uses three verdict labels. \emph{Supported} means that time-valid evidence corroborates the claim, \emph{refuted} means that it contradicts the claim, and \emph{insufficient evidence} means that the available material cannot settle the claim. Selective abstention is a separate model decision: the system may defer any prediction when its calibrated confidence is too low. This separation permits independent measurement of evidential sufficiency and decision reliability.

Recent systems address parts of this setting. MOCHEG joins multimodal retrieval and verification \cite{yao2023mocheg}; RED-DOT reranks web evidence by relevance \cite{papadopoulos2025reddot}; and COVE reconstructs an image's original context before assessing its caption \cite{tonglet2025cove}. ChronoFact models temporal relations within claims and evidence \cite{barik2025chronofact}, while work on conflicting evidence shows that hard credibility filtering can remove useful counter-evidence \cite{ge2025conflicting}. None of these components alone establishes that retrieved evidence was available at the claim-specific decision time or that repeated documents represent independent sources.

This study asks whether decision-time-valid evidence from independent source clusters improves future-event fact-checking and selective risk relative to matched document-level multimodal RAG. Each case is an English claim about a United States event, paired with an image and a decision cutoff. \systemname{} admits only evidence with a documented pre-cutoff observation, groups dependent documents, represents the resulting trail in a typed graph, and estimates evidence sufficiency before veracity. Source dependence is not treated as source credibility: the primary model limits duplicate influence without suppressing evidence because of a publisher-level trust score.

The study makes three contributions:
\begin{itemize}
    \item a decision-time construction protocol for \datasetname{} that records timestamp confidence, excludes unobserved pre-cutoff evidence from the primary analysis, and isolates events, near-duplicate claims, and reused images by split;
    \item a source-dependence-aware framework that bounds each evidence cluster's aggregate mass, separates audited sufficiency supervision from inherited verdicts, and applies calibrated selective deferral; and
    \item a falsifiable evaluation using future events and counterfactual evidence interventions for post-cutoff leakage, duplicate inflation, early-source removal, and dependent versus independent conflicts.
\end{itemize}

